\section{Supervised Learning on PCIB Signals}
\label{sec:stacking}

While the theory-guided PCIB framework provides interpretable hallucination signals, 
we can further improve performance by treating these signals as features in supervised 
machine learning models. This hybrid approach combines the theoretical foundations of PCIB 
with the empirical optimization of data-driven methods.

\subsection{Feature Engineering}
We extract and engineer features from PCIB signals:

\begin{itemize}
  \item \textbf{Base Signals}: Uptake KL-divergence, Stress JS-divergence, Conflict JS-divergence, and composite score
  \item \textbf{Aggregations}: Mean, maximum, minimum, and standard deviation across all claims
  \item \textbf{Interactions}: Pairwise and three-way interactions between signals (e.g., $\text{uptake} \times \text{stress}$)
  \item \textbf{Ratios}: Relative signal strengths (e.g., $\text{uptake}/\text{stress}$)
  \item \textbf{Polynomials}: Squared terms to capture non-linear relationships
  \item \textbf{Composite Features}: Products of composite score with individual signals
\end{itemize}

This feature engineering process expands the 11 base PCIB features to over 30 engineered features, 
capturing complex non-linear relationships and interactions between signals.

\subsection{Model Architecture}
We train multiple supervised learning models on the engineered PCIB features:

\begin{itemize}
  \item \textbf{Random Forest}: Ensemble of 500 decision trees with balanced class weights
  \item \textbf{Gradient Boosting}: Sequential tree-based learner with learning rate 0.01
  \item \textbf{Neural Network}: Multi-layer perceptron with hidden layers [100, 50]
  \item \textbf{Support Vector Machine}: RBF kernel with isotonic calibration
\end{itemize}

All models use 5-fold stratified cross-validation for hyperparameter tuning and evaluation. 
We optimize for AUROC during grid search and apply isotonic calibration where appropriate.

\subsection{Ensemble Methods}
We explore two ensemble strategies:

\begin{enumerate}
  \item \textbf{Simple Average}: Arithmetic mean of all model predictions
  \item \textbf{Optimized Weighted}: Grid search over weight combinations to maximize AUROC
\end{enumerate}

The optimized weighted ensemble searches through weight combinations and selects 
the configuration that maximizes validation AUROC, balancing the strengths of different model types.

\subsection{Results}
Table~\ref{tab:stacking_performance} presents the performance of each approach. 
The baseline PCIB composite score achieves AUROC = 0.8017, while 
the best stacking method (Optimized Weighted Ensemble) achieves AUROC = 0.8543, 
a relative improvement of 6.56\%. 

Figure~\ref{fig:stacking_comparison} visualizes the performance gains across different methods. 
All supervised learning approaches outperform the baseline, with ensemble methods achieving 
the strongest performance. This demonstrates that while PCIB signals are theoretically motivated 
and interpretable, data-driven optimization can further enhance their effectiveness.

\subsection{Analysis}
The success of stacking approaches reveals several insights:

\begin{itemize}
  \item \textbf{Non-linear Patterns}: Gradient Boosting and Neural Networks capture non-linear relationships between PCIB signals
  \item \textbf{Feature Importance}: Tree-based models reveal that composite score and uptake KL-divergence are most predictive
  \item \textbf{Ensemble Benefits}: Combining diverse models (tree-based, kernel-based, neural) captures complementary patterns
  \item \textbf{Calibration Matters}: Isotonic calibration improves probability estimates, especially for Random Forest and SVM
\end{itemize}

Importantly, this approach maintains the interpretability of PCIB signals while achieving 
competitive performance with state-of-the-art hallucination detection systems. The learned 
feature weights can be analyzed to understand which signal combinations are most indicative 
of hallucinations in different contexts.